% ============================================================================
\chapter{Construire une API REST en PHP (et l'authentifier avec un JWT)}
\label{ch-php-api}
% ============================================================================

\begin{objectif}
\begin{itemize}
  \item Concevoir une API coh\'erente~: ressources, verbes, codes, pagination, erreurs, versionnage.
  \item \'Ecrire en PHP pur un mini-framework~: routeur, middlewares, validation, s\'erialisation.
  \item Authentifier par jeton (JWT) et par cookie, et savoir quand employer l'un ou l'autre.
  \item Livrer une documentation exploitable (OpenAPI) et un client de test.
  \item Prot\'eger le point d'entr\'ee~: CORS, d\'ebit, tailles, quotas, idempotence.
\end{itemize}
\end{objectif}

\section{D'abord le contrat (10 minutes qui en sauvent 100)}

\begin{tabz}{lXX}
  \headrow \thd{Question} \& \thd{R\'eponse par d\'efaut} \& \thd{Et si tu fais autrement} \\
  \midrule
  Mod\`ele d'URL \& \verb|/api/v1/ressource/{id}/sous-ressource| \& \verb|/articles/lire?id=..| devient invendable. \\
  Verbes \& GET lire, POST cr\'eer, PUT remplacer, PATCH modifier partiellement, DELETE supprimer \& un POST partout cache la s\'emantique et casse le cache. \\
  Codes \& 200, 201, 204, 400, 401, 403, 404, 409, 422, 429, 500 \& 200 avec \verb|{erreur}| force le client \`a lire le corps pour savoir. \\
  Erreurs \& RFC 9457 \verb|application/problem+json| \& un format maison par endpoint~: enfer. \\
  Pagination \& \verb|?page=\&limit=| + en-t\^ete \verb|Link|, ou curseur pour les gros volumes \& sans limite maximale, une requ\^ete tue la base. \\
  Dates \& ISO~8601 UTC (\verb|2026-01-31T10:00:00Z|) \& \verb|31/01/2026|~: ambiguity garantie. \\
  Ids \& opaque (\verb|string| ou UUIDv7) ; jamais auto-incr\'ement\'e expos\'e seul \& \'enum\'erable : les voisins lisent ton catalogue. \\
  Version \& dans l'URL \verb|/v1/| ou un en-t\^ete accept ; une seule r\`egle, la m\^eme partout \& \verb|?v=1| partout : le chaos. \\
\end{tabz}

\begin{codesql}
-- Exemple de contrat d'erreur unique (RFC 9457) :
{
  "type": "https://api.nour.dev/problemes/validation",
  "title": "Donnees invalides",
  "status": 422,
  "detail": "3 champs rejettent la validation.",
  "instance": "/api/v1/articles",
  "trace": "a91f3c",
  "errors": [
    { "champ": "titre", "code": "trop_court", "message": "Au moins 3 caracteres." },
    { "champ": "tags[2]", "code": "trop_long", "message": "30 caracteres maximum." }
  ]
}
\end{codesql}

\section{Le mini-framework~: 150 lignes qui expliquent tout}

\begin{codephp}
<?php
// src/Noyau/App.php --- un noyau d'API, sans dependance, a comprendre ligne par ligne
declare(strict_types=1);
namespace App\Noyau;

use JsonException;

final class Requete
{
    /** @param array<string,mixed> $params @param array<string,mixed> $corps */
    public function __construct(
        public readonly string $methode,
        public readonly string $chemin,
        public readonly array $params = [],
        public readonly array $corps = [],
        public readonly array $query = [],
        public readonly array $entetes = [],
        public ?object $utilisateur = null,
    ) {}
    public function entete(string $nom): ?string { return $this->entetes[strtolower($nom)] ?? null; }
}

final class Reponse
{
    /** @param array<string,string> $entetes */
    public function __construct(
        public readonly int $code = 200,
        public readonly mixed $corps = null,
        public readonly array $entetes = [],
        public readonly ?string $brut = null,
    ) {}
    public static function json(mixed $donnees, int $code = 200, array $entetes = []): self
    {
        return new self($code, $donnees, ['content-type' => 'application/json; charset=utf-8'] + $entetes);
    }
    public static function probleme(int $code, string $titre, string $detail = '', array $extra = []): self
    {
        return self::json(array_merge([
            'type' => 'https://api.nour.dev/problemes/' . slug($titre),
            'title' => $titre, 'status' => $code, 'detail' => $detail,
        ], $extra), $code, ['content-type' => 'application/problem+json; charset=utf-8']);
    }
    public static function vide(int $code = 204): self { return new self($code); }
}

function slug(string $s): string
{
    return strtolower(preg_replace('/[^a-z0-9]+/i', '-', iconv('UTF-8', 'ASCII//TRANSLIT', $s) ?: $s) ?? '') ?? 'erreur';
}

final class App
{
    /** @var array<int, array{methode:string, motif:string, gestionnaire:callable, middlewares:array}> */
    private array $routes = [];
    /** @var callable[] */ private array $middlewares = [];

    public function get(string $chemin, callable $h): void { $this->ajouter('GET', $chemin, $h); }
    public function post(string $chemin, callable $h): void { $this->ajouter('POST', $chemin, $h); }
    public function patch(string $chemin, callable $h): void { $this->ajouter('PATCH', $chemin, $h); }
    public function delete(string $chemin, callable $h): void { $this->ajouter('DELETE', $chemin, $h); }

    public function middleware(callable $m): void { $this->middlewares[] = $m; }

    private function ajouter(string $methode, string $chemin, callable $h): void
    {
        $motif = preg_replace_callback('/\{(\w+)(?::([^}]+))?\}/', static function (array $m): string {
            return sprintf('(?<%s>%s)', $m[1], $m[2] ?? '[^/]+');
        }, $chemin);
        $this->routes[] = ['methode' => $methode, 'motif' => '#^' . $motif . '$#', 'gestionnaire' => $h, 'middlewares' => []];
    }

    public function traiter(): void
    {
        $methode = $_SERVER['REQUEST_METHOD'] ?? 'GET';
        $chemin = parse_url($_SERVER['REQUEST_URI'] ?? '/', PHP_URL_PATH) ?: '/';
        $entetes = [];
        foreach ($_SERVER as $k => $v) {
            if (str_starts_with((string) $k, 'HTTP_')) {
                $entetes[strtolower(str_replace('_', '-', substr((string) $k, 5)))] = (string) $v;
            }
        }
        $corps = $this->lireCorps($entetes);

        foreach ($this->routes as $r) {
            if ($r['methode'] !== $methode || !preg_match($r['motif'], $chemin, $m)) { continue; }
            $params = array_filter($m, 'is_string', ARRAY_FILTER_USE_KEY);
            $req = new Requete($methode, $chemin, $params, $corps, $_GET, $entetes);
            $chain = array_reduce(
                array_reverse([...$this->middlewares, static fn (Requete $q) => ($r['gestionnaire'])($q)]),
                static fn (callable $suivant, callable $mw): callable => static fn (Requete $q): Reponse => $mw($q, $suivant),
                static fn (Requete $q): Reponse => new Reponse(500)
            );
            $this->envoyer($chain($req));
            return;
        }
        $this->envoyer(Reponse::probleme(404, 'Ressource introuvable', $methode . ' ' . $chemin));
    }

    /** @return array<string,mixed> */
    private function lireCorps(array $entetes): array
    {
        $type = $entetes['content-type'] ?? '';
        $brut = file_get_contents('php://input') ?: '';
        if (str_contains($type, 'json')) {
            if ($brut === '') { return []; }
            try { return (array) json_decode($brut, true, 32, JSON_THROW_ON_ERROR); }
            catch (JsonException $e) { throw new ErreurRequete('JSON invalide : ' . $e->getMessage(), 400); }
        }
        if (str_contains($type, 'x-www-form-urlencoded')) { parse_str($brut, $out); return $out; }
        return $_POST;
    }

    private function envoyer(Reponse $r): void
    {
        http_response_code($r->code);
        foreach ($r->entetes as $nom => $valeur) { header($nom . ': ' . $valeur, true); }
        if ($r->corps !== null) {
            echo json_encode($r->corps, JSON_UNESCAPED_UNICODE | JSON_UNESCAPED_SLASHES | JSON_THROW_ON_ERROR);
        }
    }
}

final class ErreurRequete extends \RuntimeException { public function __construct(string $m, public readonly int $code = 400) { parent::__construct($m); } }
\end{codephp}

\begin{codephp}
<?php
// public/index.php --- branchement : middlewares, gestion d'erreur, routes
declare(strict_types=1);
require __DIR__ . '/../vendor/autoload.php';

use App\Noyau\{App, Requete, Reponse, ErreurRequete};

$app = new App();
$app->middleware(function (Requete $r, callable $suite): Reponse {
    header('x-request-id: ' . bin2hex(random_bytes(6)));
    header('content-type: application/json; charset=utf-8');
    return $suite($r);
});
$app->middleware(App\Securite::cors(...));                 // cf. chapitre securite
$app->middleware(App\Securite::limiteurDebit(120, 60));    // 120 req/min/IP (apcu)
$app->middleware(App\Auth::porte('bearer'));                // remplit $r->utilisateur si jeton valide

$articles = new App\Depot\DepotArticles(pdo());
$valider = new App\Validation\Valideur();

$app->get('/api/v1/sante', fn () => Reponse::json(['ok' => true, 'php' => PHP_VERSION, 'base' => pdo()->query('SELECT 1')->fetchColumn() ? 'ok' : 'ko']));

$app->get('/api/v1/articles', function (Requete $r) use ($articles, $valider) {
    $q = $valider->verifier($r->query, [
        'page' => 'int|min:1', 'limite' => 'int|min:1|max:100', 'recherche' => 'string|max:120',
    ]);
    return Reponse::json($articles->pageListe($q['page'] ?? 1, $q['limite'] ?? 20, $q['recherche'] ?? ''));
});

$app->post('/api/v1/articles', function (Requete $r) use ($articles, $valider) {
    $donnees = $valider->verifier($r->corps, [
        'titre' => 'requis|string|min:3|max:160',
        'corps' => 'requis|string|min:20',
        'tags' => 'array|max:12',
    ], [
        'tags.*' => 'string|min:1|max:30',
    ]);
    if ($r->utilisateur === null) { return Reponse::probleme(401, 'Connexion requise'); }
    $cree = $articles->creer($donnees, $r->utilisateur);
    return Reponse::json($cree, 201)->avecEntete('location', '/api/v1/articles/' . $cree['id']);
});

$app->patch('/api/v1/articles/{id:\d+}', function (Requete $r) use ($articles) {
    $maj = $articles->modifier((int) $r->params['id'], $r->corps, $r->utilisateur);
    return $maj === null ? Reponse::probleme(404, 'Introuvable') : Reponse::json($maj);
});

$app->delete('/api/v1/articles/{id:\d+}', fn (Requete $r) => $articles->supprimer((int) $r->params['id'], $r->utilisateur)
    ? Reponse::vide() : Reponse::probleme(404, 'Introuvable'));

try {
    $app->traiter();
} catch (ErreurRequete $e) {
    (function (Reponse $r): void { http_response_code($r->code); foreach ($r->entetes as $n => $v) { header("$n: $v"); } echo json_encode($r->corps); })
    (Reponse::probleme($e->code, 'Requete invalide', $e->getMessage()));
} catch (Throwable $e) {
    error_log((string) $e);
    (static function (): void { http_response_code(500); header('content-type: application/problem+json');
        echo json_encode(['title' => 'Erreur interne', 'status' => 500, 'trace' => substr(hash('crc32b', (string) $e), 0, 8)]); })();
}
exit;
\end{codephp}

\subsection{Un validateur minimal}

\begin{codephp}
<?php
declare(strict_types=1);
namespace App\Validation;

final class Valideur
{
    /** @param array<string,string> $regles  @return array<string,mixed> @throws ErreurValidation */
    public function verifier(array $donnees, array $regles, array $reglesNichees = []): array
    {
        $erreurs = []; $out = [];
        foreach ($regles as $champ => $chaine) {
            $r\`egles = explode('|', $chaine);
            $valeur = $donnees[$champ] ?? null;
            $obligatoire = in_array('requis', $r\`egles, true);
            if ($valeur === null || $valeur === '') {
                if ($obligatoire) { $erreurs[$champ][] = 'Champ obligatoire.'; }
                continue;
            }
            foreach ($r\`egles as $r) {
                [$nom, $arg] = array_pad(explode(':', $r, 2), 2, null);
                $message = match ($nom) {
                    'requis' => null,
                    'string' => is_string($valeur) ? null : 'Doit etre une chaine.',
                    'int' => filter_var($valeur, FILTER_VALIDATE_INT) === false ? 'Nombre entier attendu.' : null,
                    'bool' => null,
                    'array' => is_array($valeur) ? null : 'Liste attendue.',
                    'email' => filter_var($valeur, FILTER_VALIDATE_EMAIL) === false ? 'Email invalide.' : null,
                    'url' => filter_var($valeur, FILTER_VALIDATE_URL) === false ? 'URL invalide.' : null,
                    'min' => mb_strlen((string) $valeur) < (int) $arg ? sprintf('Minimum %s caract\`eres.', $arg) : null,
                    'max' => mb_strlen((string) $valeur) > (int) $arg ? sprintf('Maximum %s caract\`eres.', $arg) : null,
                    'max_count' => is_array($valeur) && count($valeur) > (int) $arg ? sprintf('Maximum %s elements.', $arg) : null,
                    'in' => !in_array((string) $valeur, explode(',', (string) $arg), true) ? 'Valeur non autoris\'ee.' : null,
                    default => throw new InvalidArgumentException("regle inconnue : {$r}"),
                };
                if ($message !== null) { $erreurs[$champ][] = $message; break; }
            }
            $out[$champ] = match (true) {
                in_array('int', $r\`egles, true) => (int) $valeur,
                in_array('bool', $r\`egles, true) => filter_var($valeur, FILTER_VALIDATE_BOOL),
                in_array('string', $r\`egles, true) => trim((string) $valeur),
                default => $valeur,
            };
        }
        if ($erreurs !== []) { throw new ErreurValidation($erreurs); }
        return $out;
    }
}

final class ErreurValidation extends \RuntimeException
{
    /** @param array<string,list<string>> $erreurs */
    public function __construct(public readonly array $erreurs)
    {
        parent::__construct('Validation echouee');
    }
    public function versProblem(): array
    {
        return ['errors' => array_map(
            static fn (array $l): array => ['message' => $l[0]],
            $this->erreurs
        )];
    }
}
\end{codephp}

\section{JWT~: le fabriquer, le v\'erifier, et savoir quand s'en passer}

\begin{codephp}
<?php
declare(strict_types=1);
namespace App\Auth;

/**
 * JWT minimal (HS256) en 40 lignes : tu comprends enfin ce que font les bibliotheques.
 * En production, prends firebase/php-jwt ou lcobucci/jwt -- et surtout la config qui va avec.
 */
final class Jeton
{
    public function __construct(private readonly string $secret, private readonly int $duree = 900) {}

    /** @param array<string,mixed> $revendications */
    public function emettre(int $utilisateurId, array $revendications = []): string
    {
        $entete = ['typ' => 'JWT', 'alg' => 'HS256'];
        $maintenant = time();
        $corps = [
            'iss' => 'https://api.nour.dev',
            'aud' => 'nour-web',
            'sub' => (string) $utilisateurId,
            'iat' => $maintenant,
            'nbf' => $maintenant,
            'exp' => $maintenant + $this->duree,
            'jti' => bin2hex(random_bytes(8)),
        ] + $revendications;
        $segments = [
            self::b64($entete),
            self::b64($corps),
        ];
        $signature = hash_hmac('sha256', implode('.', $segments), $this->secret, true);
        $segments[] = self::b64($signature, brut: true);
        return implode('.', $segments);
    }

    /** @return array<string,mixed> @throws ErreurJeton */
    public function verifier(string $jeton): array
    {
        $parties = explode('.', $jeton);
        if (count($parties) !== 3) { throw new ErreurJeton('Format invalide'); }
        [$h, $c, $s] = $parties;
        $attendu = hash_hmac('sha256', $h . '.' . $c, $this->secret, true);
        $recu = base64_decode(strtr($s, '-_', '+/'), true);
        if ($recu === false || !hash_equals(hash('sha256', $attendu), hash('sha256', $recu))) {
            throw new ErreurJeton('Signature invalide');
        }
        $entete = self::deb64($h);
        // Contre l'attaque "alg: none" : algorithme epingle
        if (($entete['alg'] ?? null) !== 'HS256') { throw new ErreurJeton('Algorithme non autorise'); }
        $corps = self::deb64($c);
        $now = time();
        foreach (['nbf' => '<', 'iat' => '<', 'exp' => '>'] as $cle => $sens) {
            if (isset($corps[$cle]) && ($sens === '<' ? $now < (int) $corps[$cle] : $now > (int) $corps[$cle])) {
                throw new ErreurJeton("Jeton non valide pour {$cle}");
            }
        }
        foreach (['iss' => 'https://api.nour.dev', 'aud' => 'nour-web'] as $cle => $valeur) {
            if (($corps[$cle] ?? null) !== $valeur) { throw new ErreurJeton("Revendication {$cle} invalide"); }
        }
        return $corps;
    }

    private static function b64(array|string $donnees, bool $brut = false): string
    {
        $chaine = $brut ? $donnees : json_encode($donnees, JSON_UNESCAPED_SLASHES | JSON_THROW_ON_ERROR);
        return rtrim(strtr(base64_encode($chaine), '+/', '-_'), '=');
    }
    private static function deb64(string $s): array
    {
        $d = base64_decode(strtr($s, '-_', '+/'), true);
        if ($d === false) { throw new ErreurJeton('Base64 invalide'); }
        return (array) json_decode($d, true, 16, JSON_THROW_ON_ERROR);
    }
}
final class ErreurJeton extends \RuntimeException {}

// Le middleware : autorisation = verifier le jeton, puis CHARGER l'utilisateur en base.
// Ne fais jamais confiance aux roles stockes dans le jeton pour une decision sensible
// (le jeton vit 15 min ; un role revoque doit etre effectif immediatement).
final class Auth
{
    public static function porte(string $type): callable
    {
        return static function (\App\Noyau\Requete $r, callable $suite) use ($type): \App\Noyau\Reponse {
            $jeton = null;
            if ($type === 'bearer' && preg_match('/^Bearer\s+(\S+)$/i', (string) $r->entete('authorization'), $m)) {
                $jeton = $m[1];
            }
            if ($jeton !== null) {
                try {
                    $corps = (new Jeton($_ENV['JWT_SECRET']))->verifier($jeton);
                    $r = $r->avecUtilisateur(Utilisateurs::trouver((int) $corps['sub']));
                } catch (ErreurJeton $e) {
                    return \App\Noyau\Reponse::probleme(401, 'Jeton invalide', $e->getMessage());
                }
            }
            return $suite($r);
        };
    }
}
\end{codephp}

\begin{tabz}{llX}
  \headrow \thd{Point} \& \thd{Cookie de session} \& \thd{JWT bearer} \\
  \midrule
  R\'evoquer \& instantan\'e (on supprime la ligne) \& impossible avant expiration~: pr\'evoir une liste de r\'evocation \\
  CSRF \& \`a prot\'eger (jeton) \& quasi inexistant (le navigateur ne l'ajoute pas tout seul) \\
  Multi-sous-domaines \& possible (\verb|Domain=.nour.dev|) \\
  Mobile / API tierces \& awkward \& naturel \\
  \'Etat serveur \& requis (Redis/BDD) \& aucun (mais cache quand m\^eme) \\
  S\'ecurit\'e du stockage \& le JS ne le voit jamais (\verb|HttpOnly|) \& dans \verb|localStorage| = vol\'e par le moindre XSS \\
\end{tabz}

\begin{astuce}
Pour une application web JS + PHP~: \textbf{cookie de session} (Simple, CSRF pr\`es, r\'evocable) pour le
navigateur, \textbf{JWT} pour le mobile et les int\'egrations tierces. Les deux pointent sur le m\^eme
contr\^oleur~: le middleware produit un objet utilisateur, le reste du code ne sait pas comment il a \'et\'e
authentifi\'e.
\end{astuce}

\subsection{CORS, quotas, idempotence~: les trois finitions}

\begin{codephp}
<?php
final class Securite
{
    public static function cors(\App\Noyau\Requete $r, callable $suite): \App\Noyau\Reponse
    {
        $origines = ['http://localhost:5173', 'https://nour.dev'];
        $o = $r->entete('origin');
        $autorisee = $o !== null && in_array($o, $origines, true);
        $enTetes = static function (array $h) use ($autorisee, $o): array {
            if (!$autorisee) { return $h; }
            return $h + [
                'access-control-allow-origin' => $o,
                'vary' => 'origin',
                'access-control-allow-credentials' => 'true',
                'access-control-max-age' => '600',
            ];
        };
        if ($r->methode === 'OPTIONS') {
            return new \App\Noyau\Reponse(204, null, $enTetes([
                'access-control-allow-methods' => 'GET,POST,PUT,PATCH,DELETE,OPTIONS',
                'access-control-allow-headers' => 'content-type,authorization,idempotency-key',
            ]));
        }
        return $suite($r);
    }

    public static function limiteurDebit(int $max, int $fenetre): callable
    {
        return static function (\App\Noyau\Requete $r, callable $suite) use ($max, $fenetre): \App\Noyau\Reponse {
            if (function_exists('apcu_add') && $r->utilisateur === null) {
                $cle = 'rl:' . ($_SERVER['REMOTE_ADDR'] ?? '?') . ':' . (int) floor(time() / $fenetre);
                if (!apcu_add($cle, 1, $fenetre)) {
                    $n = apcu_inc($cle);
                    if ($n > $max) {
                        return \App\Noyau\Reponse::probleme(429, 'Trop de requetes', "limite : {$max}/{$fenetre}s",
                            ['retry_after' => $fenetre - (time() % $fenetre)]);
                    }
                }
            }
            return $suite($r);
        };
    }
}

// Idempotence pour les POST critiques (paiement, envoi d'email) :
// une cle client -> on rejoue la meme reponse au lieu de creer un doublon.
function avecCleIdempotence(callable $creation): callable
{
    return static function (array $donnees, string $cle) use ($creation) {
        $pdo = pdo();
        $pdo->prepare('INSERT IGNORE INTO idempotence (cle, cree_le) VALUES (?, NOW())')->execute([$cle]);
        $dejaLa = $pdo->prepare('SELECT reponse FROM idempotence WHERE cle = ? AND reponse IS NOT NULL');
        $dejaLa->execute([$cle]);
        if ($deja = $dejaLa->fetchColumn()) { return json_decode((string) $deja, true, 512, JSON_THROW_ON_ERROR); }
        $reponse = $creation($donnees);
        $pdo->prepare('UPDATE idempotence SET reponse = ? WHERE cle = ?')
            ->execute([json_encode($reponse, JSON_THROW_ON_ERROR), $cle]);
        return $reponse;
    };
}
\end{codephp}

\section{Documenter et tester~: OpenAPI + un client}

\begin{codeyml}
# openapi.yaml (extrait) --- a editer ou a generer depuis les attributs + zircote/swagger-php
openapi: 3.1.0
info: { title: "API Nour", version: "1.0.0", description: "Articles, tags, auth." }
servers: [{ url: "https://api.nour.dev/v1" }]
paths:
  /articles:
    get:
      summary: Liste paginee
      parameters:
        - { name: page, in: query, schema: { type: integer, minimum: 1, default: 1 } }
        - { name: limite, in: query, schema: { type: integer, maximum: 100, default: 20 } }
      responses:
        "200":
          description: OK
          content:
            application/json:
              schema: { $ref: "#/components/schemas/PageArticles" }
        "429": { $ref: "#/components/responses/TropDeRequetes" }
components:
  schemas:
    Article:
      type: object
      required: [id, titre, statut]
      properties:
        id: { type: integer, format: int64, readOnly: true }
        titre: { type: string, minLength: 3, maxLength: 160 }
        statut: { type: string, enum: [brouillon, publie, archive] }
        publie_le: { type: string, format: date-time, nullable: true }
\end{codeyml}

\begin{terminal}
# Generer la spec depuis le code : composer require --dev zircote/swagger-php
vendor/bin/openapi src -o public/openapi.json --format yaml
# Servir la doc : https://github.com/scalar/scalar (Scalar, le meilleur lecteur moderne)
npx @scalar/api-reference --url /openapi.json
# Valider la spec + generer des exemples de contrats :
npx @redocly/cli lint openapi.yaml
# Client de test dans VS Code : api.http
\end{terminal}

\begin{codeplain}
### api.http --- extension "REST Client" (humao.rest-client)
@base = http://127.0.0.1:8000/api/v1

GET {{base}}/sante

###
# @name login
POST {{base}/auth/connexion
content-type: application/json

{ "email": "nour@example.com", "mot_de_passe": "un mot de passe costaud" }

###
@token = {{login.response.body.jeton}}

GET {{base}}/articles?page=1
authorization: Bearer {{token}}
x-request-id: 01HTQ...

###
POST {{base}}/articles
authorization: Bearer {{token}}
idempotency-key: 01HTQ6M2K9
content-type: application/json

{ "titre": "Apprendre PHP vraiment", "corps": "Un article long et honnete...", "tags": ["php"] }
\end{codeplain}

\begin{projet}[API compl\`ete]
Construis l'API de ton blog~: authentification (inscription, connexion, refresh, d\'econnexion,
r\'einitialisation), articles CRUD avec upload d'image de couverture, tags, commentaires
(mod\'erables), recherche, pagination par curseur, quota par utilisateur, OpenAPI servi sur
\verb|/docs|, 40 tests d'int\'egration (HTTP r\'eel + SQLite en m\'emoire), et un client JS g\'en\'er\'e
\`a partir de la spec (\link{openapi-typescript-codegen}{https://github.com/feross/openapi-typescript-codegen}
ou \link{oazapfts}{https://github.com/oazapfts/oazapfts}). C'est le projet qui, sur un CV, se
distingue des dizaines de \og ~to-do app~ \fg{}.
\end{projet}

\begin{ressources}
\begin{itemize}
  \item RFC 9457 (Problem Details) --- \url{https://datatracker.ietf.org/doc/html/rfc9457}
  \item OpenAPI~3.1 --- \url{https://spec.openapis.org/oas/latest.html} --- Scalar \url{https://scalar.com/}
  \item \link{Firebase~PHP~JWT}{https://github.com/firebase/php-jwt} --- \link{LCobucci~JWT}{https://github.com/lcobucci/jwt}
  \item API design guide (Microsoft) --- \url{https://github.com/microsoft/api-guidelines}
  \item \link{Zircote Swagger-PHP (attributs -> OpenAPI)}{https://github.com/zircote/swagger-php}
\end{itemize}
\end{ressources}

\index{API REST}\index{JWT}\index{OpenAPI}\index{CORS}
